% Reference entries in the Applied Energy (Elsevier numbered) style. They are numbered in the order of first citation
% by tools/order_refs.py, which writes sections/references_ordered.tex. Entries marked [verify] should be checked
% against the publisher's page before submission.
\bibitem{iea2025} International Energy Agency. Energy and AI. Paris: IEA; 2025.
\bibitem{shehabi2024} Shehabi A, Smith SJ, Hubbard A, et al. 2024 United States data center energy usage report. Berkeley, CA: Lawrence Berkeley National Laboratory; 2024.
\bibitem{ndrc2024dc} National Development and Reform Commission, Ministry of Industry and Information Technology, National Energy Administration, National Data Administration. Special action plan for the green and low-carbon development of data centers. Beijing; 2024 [in Chinese].
\bibitem{norris2025} Norris TH, Profeta T, Pati\~no-Echeverri D, Cowie-Haskell A. Rethinking load growth: assessing the potential for integration of large flexible loads in US power systems. Durham, NC: Nicholas Institute for Energy, Environment \& Sustainability, Duke University; 2025.
\bibitem{chen2026defer} Chen Y, Zheng X. To defer or to shift? The role of AI data center flexibility on grid interconnection. arXiv:2604.05376; 2026.
\bibitem{williams2026} Williams C, Colangelo P, Coskun AK, et al. Power-flexible AI data centers: a new paradigm for grid-responsive compute. arXiv:2606.25098; 2026.
\bibitem{texas2025sb6} Texas Legislature. Senate Bill 6, 89th Legislature: interconnection and curtailment of large loads. Austin, TX; 2025.
\bibitem{ndrc2023dsm} National Development and Reform Commission, et al. Measures for the administration of power demand-side management (2023 edition). Beijing; 2023 [in Chinese].
\bibitem{jiangsu2024dr} Jiangsu Provincial Development and Reform Commission. Implementation rules for power demand response in Jiangsu Province. Nanjing; 2024 [in Chinese].
\bibitem{colangelo2025} Colangelo P, Coskun AK, Megrue J, et al. Turning AI data centers into grid-interactive assets: results from a field demonstration in Phoenix, Arizona. arXiv:2507.00909; 2025.
\bibitem{liu2013} Liu Z, Wierman A, Chen Y, Razon B, Chen N. Data center demand response: avoiding the coincident peak via workload shifting and local generation. Perform Eval 2013;70(10):770--91.
\bibitem{wierman2014} Wierman A, Liu Z, Liu I, Mohsenian-Rad H. Opportunities and challenges for data center demand response. In: Proceedings of the International Green Computing Conference (IGCC); 2014.
\bibitem{chen2020review} Chen M, Gao C, Song M, et al. Internet data centers participating in demand response: a comprehensive review. Renew Sustain Energy Rev 2020;117:109466.
\bibitem{radovanovic2023} Radovanovi\'c A, Koningstein R, Schneider I, et al. Carbon-aware computing for datacenters. IEEE Trans Power Syst 2023;38(2):1270--80.
\bibitem{acun2026} Acun F, Hankendi C, Levine E, Reynolds H, Bardwick J, Coskun AK. Investigating power consumption flexibility of AI data centers for demand response participation. In: Proceedings of the ACM International Conference on Future and Sustainable Energy Systems (e-Energy); 2026.
\bibitem{patel2024} Patel P, Choukse E, Zhang C, et al. Characterizing power management opportunities for LLMs in the cloud. In: Proceedings of the ACM International Conference on Architectural Support for Programming Languages and Operating Systems (ASPLOS); 2024, p. 207--22.
\bibitem{du2026} Du B, Fan X, Du E, Chen L, Han J, Hou Q, et al. From tokens to energy flexibility: quantization-enabled demand response for data centers with LLM inference workloads. arXiv:2606.18851; 2026.
\bibitem{google2025} Google. Demand response agreements with Indiana Michigan Power and the Tennessee Valley Authority for machine-learning workloads. Company announcement, 4 August 2025.
\bibitem{li2026scalar} Li M. Beyond scalar flexibility: from eligible AI workloads to dependable load relief. arXiv:2609.05406; 2026.
\bibitem{choukse2025} Choukse E, et al. Power stabilization for AI training datacenters. arXiv:2508.14318; 2025.
\bibitem{han2024} Han O, Ding T, Yang M, et al. A novel 4-level joint optimal dispatch for demand response of data centers with district autonomy realization. Appl Energy 2024;358:122590.
\bibitem{lin2025} Lin J, Han J, Wang Y, et al. Planning and operation optimization of a computing-power integrated energy system for low-carbon data centers considering exergy efficiency. J Shanghai Jiao Tong Univ 2025;59(9):1327--37 [in Chinese].
\bibitem{fan2026} Fan H, Chen S, Luo H, et al. Optimal operation strategy of data centers considering cooling--heating--power multi-energy complementarity and spatio-temporally adjustable loads. South Power Syst Technol 2026;20(5):28--36 [in Chinese].
\bibitem{zhang2026bilevel} Zhang S, Ding S, Li Y, Wei M, Tan Q. A bi-level planning model of integrated energy capacity configuration for data center cluster considering computing-electricity coupling operation and supply-demand uncertainties. Appl Energy 2026;425:128499.
\bibitem{huang2020} Huang P, Copertaro B, Zhang X, et al. A review of data centers as prosumers in district energy systems: renewable energy integration and waste heat reuse for district heating. Appl Energy 2020;258:114109.
\bibitem{wahlroos2018} Wahlroos M, P\"arssinen M, Rinne S, et al. Future views on waste heat utilization -- case of data centers in Northern Europe. Renew Sustain Energy Rev 2018;82(2):1749--64.
\bibitem{evans2022} Evans MP, Tindemans SH, Angeli D. Flexibility framework with recovery guarantees for aggregated energy storage devices. IEEE Trans Smart Grid 2022;13(5):3519--31.
\bibitem{hao2015} Hao H, Sanandaji BM, Poolla K, Vincent TL. Aggregate flexibility of thermostatically controlled loads. IEEE Trans Power Syst 2015;30(1):189--98.
\bibitem{wrinch2012} Wrinch M, Dennis G, El-Fouly THM, Wong S. Demand response implementation for improved system efficiency in remote communities. In: Proceedings of the IEEE Electrical Power and Energy Conference (EPEC); 2012.
\bibitem{chao2011} Chao H. Demand response in wholesale electricity markets: the choice of customer baseline. J Regul Econ 2011;39(1):68--88.
\bibitem{schildbach2014} Schildbach G, Fagiano L, Frei C, Morari M. The scenario approach for stochastic model predictive control with bounds on closed-loop constraint violations. Automatica 2014;50(12):3009--18.
\bibitem{zeng2013} Zeng B, Zhao L. Solving two-stage robust optimization problems using a column-and-constraint generation method. Oper Res Lett 2013;41(5):457--61.
\bibitem{esfahani2018} Mohajerin Esfahani P, Kuhn D. Data-driven distributionally robust optimization using the Wasserstein metric: performance guarantees and tractable reformulations. Math Program 2018;171(1--2):115--66.
\bibitem{boyd2011} Boyd S, Parikh N, Chu E, Peleato B, Eckstein J. Distributed optimization and statistical learning via the alternating direction method of multipliers. Found Trends Mach Learn 2011;3(1):1--122.
\bibitem{stojkovic2025} Stojkovic J, Zhang C, Goiri \'I, Torrellas J, Choukse E. DynamoLLM: designing LLM inference clusters for performance and energy efficiency. In: Proceedings of the IEEE International Symposium on High-Performance Computer Architecture (HPCA); 2025.
\bibitem{weng2022} Weng Q, Xiao W, Yu Y, et al. MLaaS in the wild: workload analysis and scheduling in large-scale heterogeneous GPU clusters. In: Proceedings of the 19th USENIX Symposium on Networked Systems Design and Implementation (NSDI); 2022.
\bibitem{fan2007} Fan X, Weber W-D, Barroso LA. Power provisioning for a warehouse-sized computer. In: Proceedings of the 34th Annual International Symposium on Computer Architecture (ISCA); 2007, p. 13--23.
\bibitem{hu2024acme} Hu Q, Ye Z, Wang Z, et al. Characterization of large language model development in the datacenter. In: Proceedings of the 21st USENIX Symposium on Networked Systems Design and Implementation (NSDI); 2024.
\bibitem{ashrae2021} ASHRAE Technical Committee 9.9. Thermal guidelines for data processing environments. 5th ed. Atlanta, GA: ASHRAE; 2021.
\bibitem{ndrc2025green} National Development and Reform Commission, National Energy Administration. Notice on promoting green-power direct connection in an orderly manner (Fagai Nengyuan [2025] No. 650). Beijing; 2025 [in Chinese].
\bibitem{hersbach2020} Hersbach H, Bell B, Berrisford P, et al. The ERA5 global reanalysis. Q J R Meteorol Soc 2020;146(730):1999--2049.
\bibitem{gaertner2020} Gaertner E, Rinker J, Sethuraman L, et al. Definition of the IEA 15-megawatt offshore reference wind turbine. Golden, CO: National Renewable Energy Laboratory; 2020. Report NREL/TP-5000-75698.
\bibitem{meta2018odense} Meta. Odense data center: heat recovery to the Fjernvarme Fyn district heating network. Company information; 2018.
\bibitem{ndrc2016demand} General Office of the National Development and Reform Commission. Notice on improving the implementation of the basic electricity price for two-part tariff users (Fagai Ban Jiage [2016] No. 1583). Beijing; 2016 [in Chinese].
\bibitem{mee2024ef} Ministry of Ecology and Environment, National Bureau of Statistics. Announcement on the 2022 carbon dioxide emission factors of electricity. Beijing; 2024 [in Chinese].
\bibitem{rockafellar2000} Rockafellar RT, Uryasev S. Optimization of conditional value-at-risk. J Risk 2000;2(3):21--41.
\bibitem{latif2024} Latif I, Newkirk AC, Carbone MR, et al. Empirical measurements of AI training power demand on a GPU-accelerated node. arXiv:2412.08602; 2024.
\bibitem{xu2018} Xu B, Zhao J, Zheng T, Litvinov E, Kirschen DS. Factoring the cycle aging cost of batteries participating in electricity markets. IEEE Trans Power Syst 2018;33(2):2248--59.
\bibitem{dobos2014} Dobos AP. PVWatts version 5 manual. Golden, CO: National Renewable Energy Laboratory; 2014. Report NREL/TP-6A20-62641.
\bibitem{wan2025} Wan B, Han M, Sheng Y, et al. ByteCheckpoint: a unified checkpointing system for large foundation model development. In: Proceedings of the 22nd USENIX Symposium on Networked Systems Design and Implementation (NSDI); 2025.
\bibitem{wang2023gemini} Wang Z, Jia Z, Zheng S, et al. Gemini: fast failure recovery in distributed training with in-memory checkpoints. In: Proceedings of the 29th ACM Symposium on Operating Systems Principles (SOSP); 2023.
